
\documentclass{article}
\usepackage{colortbl}
\usepackage{makecell}
\usepackage{multirow}
\usepackage{supertabular}

\begin{document}

\newcounter{utterance}

\centering \large Interaction Transcript for game `codenames', experiment `abstract', episode 4 with reflexion{-}llama{-}3.1{-}8b{-}instruct.
\vspace{24pt}

{ \footnotesize  \setcounter{utterance}{1}
\setlength{\tabcolsep}{0pt}
\begin{supertabular}{c@{$\;$}|p{.15\linewidth}@{}p{.15\linewidth}p{.15\linewidth}p{.15\linewidth}p{.15\linewidth}p{.15\linewidth}}
    \# & $\;$A & \multicolumn{4}{c}{Game Master} & $\;\:$B\\
    \hline

    \theutterance \stepcounter{utterance}  
    & & & \multicolumn{2}{p{0.3\linewidth}}{
        \cellcolor[rgb]{0.9,0.9,0.9}{
            \makecell[{{p{\linewidth}}}]{
                \texttt{\tiny{[GM$|$GM]}}
                {'hidden': {'team': ['ambivalence', 'veracity'], 'innocent': ['proliferation', 'nationalization', 'subjugation', 'socialization'], 'opponent': [], 'assassin': []}, 'revealed': {'team': {'team': [], 'innocent': [], 'opponent': [], 'assassin': []}, 'opponent': {'team': [], 'innocent': [], 'opponent': [], 'assassin': []}}}
            }
        }
    }
    & & \\ \\

    \theutterance \stepcounter{utterance}  
    & & \multicolumn{4}{p{0.6\linewidth}}{
        \cellcolor[rgb]{0.9,0.9,0.9}{
            \makecell[{{p{\linewidth}}}]{
                \texttt{\tiny{[P1$\langle$GM]}}
                \texttt{You have attempted this same task before and did not succeed. These are the plans you wrote after those attempts.} \\
\texttt{Trial \#1: New plan: } \\
\\ 
\texttt{In my previous attempts, I focused on providing a single{-}word clue that was semantically related to multiple target words. However, I failed to consider the potential for my teammate to guess an opponent's word or a distractor word, which would end their turn and give the opposing team a chance to make a turn. } \\
\\ 
\texttt{Next time, I will make sure to provide a clue that is more specific and targeted towards the team's words, rather than being too broad and potentially leading to incorrect guesses. I will also consider the relationships between the words more carefully, to avoid giving away too much information that could be used by the opposing team. For example, I could provide a clue that relates to a specific aspect or characteristic of the target words, rather than their overall meaning.} \\
\\ 
\texttt{Trial \#2: New plan: } \\
\\ 
\texttt{In my previous attempts, I focused on providing a single{-}word clue that was semantically related to multiple target words. However, I failed to consider the potential for my teammate to guess an opponent's word or a distractor word, which would end their turn and give the opposing team a chance to make a turn. } \\
\\ 
\texttt{Moreover, I failed to follow the required format in my last attempt, which resulted in the game master rejecting my message. Next time, I will make sure to provide a clue that is more specific and targeted towards the team's words, rather than being too broad and potentially leading to incorrect guesses. I will also consider the relationships between the words more carefully, to avoid giving away too much information that could be used by the opposing team. } \\
\\ 
\texttt{To ensure a well{-}formed and legal response, I will pay close attention to the required format and make sure to start my clue with 'CLUE: ', followed by a new line and then 'TARGETS: ' with a comma{-}separated list of related target words.} \\
\\ 
\texttt{Trial \#3: New plan: } \\
\\ 
\texttt{In my previous attempts, I focused on providing a single{-}word clue that was semantically related to multiple target words. However, I failed to consider the potential for my teammate to guess an opponent's word or a distractor word, which would end their turn and give the opposing team a chance to make a turn. } \\
\\ 
\texttt{Moreover, I failed to follow the required format in my last attempt, which resulted in the game master rejecting my message. Next time, I will make sure to provide a clue that is more specific and targeted towards the team's words, rather than being too broad and potentially leading to incorrect guesses. I will also consider the relationships between the words more carefully, to avoid giving away too much information that could be used by the opposing team. } \\
\\ 
\texttt{To ensure a well{-}formed and legal response, I will pay close attention to the required format and make sure to start my clue with 'CLUE: ', followed by a new line and then 'TARGETS: ' with a comma{-}separated list of related target words. } \\
\\ 
\texttt{Additionally, I will make sure to provide a clue that is more closely related to the remaining team words, rather than trying to relate to all of them. For example, if I have only two team words left, I will try to find a clue that relates to both of them, rather than trying to find a clue that relates to all four team words. This will help me to provide a more targeted and effective clue that will help my teammate to guess the remaining team words.} \\
\\ 
\\ 
\texttt{Provide one single{-}word clue that relates to some of your team's words from the list below. You can choose to how many and to which words your clue relates to.} \\
\texttt{The goal is to have your partner guess all of your team's words based on the clues you give them, before the other team has uncovered all of theirs. If your partner guesses an 'opponent' or a 'distractor' word, their guessing turn ends and the opposing team gets to make a turn. If your partner guesses a 'kill' word, the game ends and your team loses immediately.} \\
\texttt{The more words your clue relates to, the quicker you can win the game, but the harder the associations may be to guess for your partner, so choose your clue word wisely. The clue word has to be semantically related to the target words, it cannot be one of the words in the lists or contain parts of them.} \\
\\ 
\texttt{Always give your single{-}word clue and your comma{-}separated list of related target words in the following format and make your answers as short as possible, never include any other text than is required in this form:} \\
\\ 
\texttt{CLUE: <WORD>} \\
\texttt{TARGETS: <TARGETS>} \\
\\ 
\texttt{Your team words are: ambivalence, veracity.} \\
\texttt{Your opponent's team words are: .} \\
\texttt{Distractor words are: proliferation, nationalization, subjugation, socialization.} \\
\texttt{The kill words are: .} \\
            }
        }
    }
    & \\ \\

    \theutterance \stepcounter{utterance}  
    & \multicolumn{4}{p{0.6\linewidth}}{
        \cellcolor[rgb]{0.9,0.9,0.9}{
            \makecell[{{p{\linewidth}}}]{
                \texttt{\tiny{[P1$\rangle$GM]}}
                \texttt{CLUE: nuance} \\
\texttt{TARGETS: ambivalence, veracity} \\
            }
        }
    }
    & & \\ \\

    \theutterance \stepcounter{utterance}  
    & & & \multicolumn{2}{p{0.3\linewidth}}{
        \cellcolor[rgb]{0.9,0.9,0.9}{
            \makecell[{{p{\linewidth}}}]{
                \texttt{\tiny{[GM$|$GM]}}
                \texttt{nuance} \\
            }
        }
    }
    & & \\ \\

    \theutterance \stepcounter{utterance}  
    & & & \multicolumn{2}{p{0.3\linewidth}}{
        \cellcolor[rgb]{0.9,0.9,0.9}{
            \makecell[{{p{\linewidth}}}]{
                \texttt{\tiny{[GM$|$GM]}}
                ['ambivalence', 'veracity']
            }
        }
    }
    & & \\ \\

    \theutterance \stepcounter{utterance}  
    & & & \multicolumn{2}{p{0.3\linewidth}}{
        \cellcolor[rgb]{0.9,0.9,0.9}{
            \makecell[{{p{\linewidth}}}]{
                \texttt{\tiny{[GM$|$GM]}}
                {'word': 'ambivalence', 'assignment': 'team'}
            }
        }
    }
    & & \\ \\

    \theutterance \stepcounter{utterance}  
    & & & \multicolumn{2}{p{0.3\linewidth}}{
        \cellcolor[rgb]{0.9,0.9,0.9}{
            \makecell[{{p{\linewidth}}}]{
                \texttt{\tiny{[GM$|$GM]}}
                {'word': 'veracity', 'assignment': 'team'}
            }
        }
    }
    & & \\ \\

    \theutterance \stepcounter{utterance}  
    & & & \multicolumn{2}{p{0.3\linewidth}}{
        \cellcolor[rgb]{0.9,0.9,0.9}{
            \makecell[{{p{\linewidth}}}]{
                \texttt{\tiny{[GM$|$GM]}}
                \texttt{opponent has won} \\
            }
        }
    }
    & & \\ \\

    \theutterance \stepcounter{utterance}  
    & \multicolumn{4}{p{0.6\linewidth}}{
        \cellcolor[rgb]{0.95,0.95,0.95}{
            \makecell[{{p{\linewidth}}}]{
                \texttt{\tiny{[P2$\rightarrow$P1]}}
                \texttt{Post{-}episode Reflexion (not sent to GM):} \\
\texttt{New plan: } \\
\\ 
\texttt{In my previous attempts, I focused on providing a single{-}word clue that was semantically related to multiple target words. However, I failed to consider the potential for my teammate to guess an opponent's word or a distractor word, which would end their turn and give the opposing team a chance to make a turn. } \\
\\ 
\texttt{Moreover, I failed to follow the required format in my last attempt, which resulted in the game master rejecting my message. Next time, I will make sure to provide a clue that is more specific and targeted towards the team's words, rather than being too broad and potentially leading to incorrect guesses. I will also consider the relationships between the words more carefully, to avoid giving away too much information that could be used by the opposing team. } \\
\\ 
\texttt{To ensure a well{-}formed and legal response, I will pay close attention to the required format and make sure to start my clue with 'CLUE: ', followed by a new line and then 'TARGETS: ' with a comma{-}separated list of related target words. } \\
\\ 
\texttt{Additionally, I will make sure to provide a clue that is more closely related to the remaining team words, rather than trying to relate to all of them. For example, if I have only two team words left, I will try to find a clue that relates to both of them, rather than trying to find a clue that relates to all four team words. This will help me to provide a more targeted and effective clue that will help my teammate to guess the remaining team words.} \\
\\ 
\texttt{I will also try to find a clue that is more closely related to the characteristics or connotations of the team words, rather than their literal meanings. For example, if the team words are related to emotions, I might provide a clue that relates to a specific emotional state or feeling. This will help me to provide a more nuanced and effective clue that will help my teammate to guess the remaining team words.} \\
            }
        }
    }
    & & \\ \\

\end{supertabular}
}

\end{document}
